\documentclass[11pt,a4paper]{article}

% Template visual Matriz Aprova. Compile com XeLaTeX ou LuaLaTeX.
\usepackage{fontspec}
\usepackage{polyglossia}
\setmainlanguage{portuguese}
\IfFontExistsTF{Aptos}{
  \setmainfont{Aptos}
  \setsansfont{Aptos}
}{
  \IfFontExistsTF{Calibri}{
    \setmainfont{Calibri}
    \setsansfont{Calibri}
  }{
    \IfFontExistsTF{TeX Gyre Heros}{
      \setmainfont{TeX Gyre Heros}
      \setsansfont{TeX Gyre Heros}
    }{
      \setmainfont{Latin Modern Sans}
      \setsansfont{Latin Modern Sans}
    }
  }
}
\usepackage[a4paper,margin=2cm,headheight=16pt]{geometry}
\usepackage{graphicx}
\usepackage{xcolor}
\usepackage{tikz}
\usepackage{tcolorbox}
\usepackage{enumitem}
\usepackage{booktabs}
\usepackage{tabularx}
\usepackage{amssymb}
\usepackage{fancyhdr}
\usepackage{titlesec}
\usepackage{hyperref}

\definecolor{matrizlime}{HTML}{CBFF4D}
\definecolor{matrizink}{HTML}{101313}
\definecolor{matrizmuted}{HTML}{596363}
\definecolor{matrizline}{HTML}{DDE3DF}
\definecolor{matrizsoft}{HTML}{F3F7EC}

\hypersetup{colorlinks=true,linkcolor=matrizink,urlcolor=matrizink}
\pagestyle{fancy}
\fancyhf{}
\lhead{\textsf{\small MATRIZ APROVA}}
\rhead{\textsf{\small APOSTILA}}
\cfoot{\textsf{\small \thepage}}
\renewcommand{\headrulewidth}{0.4pt}
\renewcommand{\headrule}{\hbox to\headwidth{\color{matrizline}\leaders\hrule height \headrulewidth\hfill}}

\titleformat{\section}{\Large\bfseries\color{matrizink}}{\thesection}{0.6em}{}
\titleformat{\subsection}{\large\bfseries\color{matrizink}}{\thesubsection}{0.6em}{}
\setlist{nosep,leftmargin=*}

\newtcolorbox{foco}{colback=matrizsoft,colframe=matrizlime,boxrule=1pt,arc=3pt,left=10pt,right=10pt,top=8pt,bottom=8pt}
\newtcolorbox{lei}{colback=matrizink,colframe=matrizink,coltext=white,boxrule=0pt,arc=3pt,left=10pt,right=10pt,top=8pt,bottom=8pt}
\newtcolorbox{alerta}{colback=white,colframe=matrizline,boxrule=0.6pt,arc=3pt,left=10pt,right=10pt,top=8pt,bottom=8pt}

% Logo usada na capa. Caminho relativo à pasta onde o .tex é compilado
% (materiais/<disciplina>/). Ajuste se a estrutura mudar.
\newcommand{\logomatriz}{../../public/matrizaprova_temaclaro.png}

\begin{document}

\begin{titlepage}
  \thispagestyle{empty}
  \begin{center}
    \includegraphics[width=0.55\linewidth]{\logomatriz}\par
  \end{center}
  \vspace{2.2cm}
  {\sffamily\fontsize{28}{32}\selectfont\bfseries Apostila — Matemática: Matemática geral para provas\par}
  \vspace{0.4cm}
  {\sffamily\large concursos gerais, carreiras militares, ENEM, escolas técnicas e vestibulares\par}
  \vfill
  \begin{foco}
    \textsf{\textbf{Foco da apostila}}\\[3pt]
    transforme o enunciado em uma representação simples, mantenha as unidades compatíveis e confira se o resultado responde ao que foi perguntado.
  \end{foco}
  \vspace{0.7cm}
  {\sffamily\small Atualizado em 16 de agosto de 2026\quad ·\quad Cebraspe, FGV, FCC, Vunesp, EsPCEx, AFA, EEAR e ENEM\par}
  \vspace{1cm}
  {\sffamily\small matrizaprova.com\par}
\end{titlepage}

\tableofcontents
\newpage

% Conteúdo gerado entra aqui.
\section*{O que cai}

Matemática de prova não é uma coleção de fórmulas isoladas. A banca apresenta uma situação, exige a tradução para uma operação e cobra atenção a unidades, sinais, ordem de grandeza e interpretação. Este material reúne os tópicos mais recorrentes para construir uma base rápida e aplicável.

O estudo deve priorizar: operações e frações; razão, proporção e regra de três; porcentagem; equações e sistemas; produtos notáveis e fatoração; funções; progressões; geometria; trigonometria; contagem e probabilidade; estatística; e matemática financeira. Em todos eles, registre os dados, escolha a relação adequada e só depois substitua os valores.

\textbf{Regra de prova:} transforme o enunciado em uma representação simples, mantenha as unidades compatíveis e confira se o resultado responde ao que foi perguntado.

\section*{Teoria essencial}

\subsection*{1. Operações, frações e ordem de cálculo}

Respeite a ordem: parênteses, colchetes e chaves; potências e raízes; multiplicações e divisões; adições e subtrações. Operações no mesmo nível são feitas da esquerda para a direita. Em \texttt{18 - 2 x 5}, primeiro vem a multiplicação: \texttt{18 - 10 = 8}.

Frações representam divisão. Para somar ou subtrair, use denominador comum. O mínimo múltiplo comum dos denominadores reduz o trabalho. Para multiplicar, multiplique numeradores e denominadores; para dividir, multiplique pela fração inversa. Assim, \texttt{2/3 + 5/6 = 4/6 + 5/6 = 9/6 = 3/2}, enquanto \texttt{(3/4) / (2/5) = (3/4) x (5/2) = 15/8}.

Em uma fração de uma quantidade, multiplique: \texttt{3/5 de 40 = (3 x 40)/5 = 24}.

Potências obedecem a \texttt{a\textasciicircum{}m x a\textasciicircum{}n = a\textasciicircum{}(m+n)}, \texttt{a\textasciicircum{}m / a\textasciicircum{}n = a\textasciicircum{}(m-n)} para \texttt{a} diferente de zero, e \texttt{(a\textasciicircum{}m)\textasciicircum{}n = a\textasciicircum{}(m x n)}. A raiz quadrada é o número não negativo cujo quadrado produz o radicando. Em provas, simplifique antes de calcular e identifique se o resultado deve ser inteiro, decimal ou fração.

\subsection*{2. Razão, proporção e regra de três}

Razão é uma comparação por quociente: a razão entre \texttt{a} e \texttt{b} é \texttt{a/b}, com \texttt{b} diferente de zero. Uma proporção é uma igualdade entre razões. Em \texttt{a/b = c/d}, o produto dos meios é igual ao produto dos extremos: \texttt{a x d = b x c}.

Na regra de três direta, as grandezas crescem ou diminuem juntas. Se 4 cadernos custam 28 reais, 7 cadernos custam \texttt{28 x 7 / 4 = 49} reais. Na regra inversa, uma grandeza aumenta enquanto a outra diminui. Se 6 trabalhadores fazem uma tarefa em 10 dias, mantendo o mesmo ritmo, 12 trabalhadores fazem em \texttt{6 x 10 / 12 = 5} dias.

Em regra de três composta, compare cada grandeza com a incógnita, marcando-a como direta ou inversa. Uma alternativa segura é montar a relação usando produtos: horas de trabalho = trabalhadores x dias x horas por dia, quando a produção é fixa. Antes de resolver, converta minutos em horas, metros em quilômetros e percentuais em números decimais quando necessário.

\subsection*{3. Porcentagem}

Porcentagem é uma razão de denominador 100. \texttt{p\%} equivale a \texttt{p/100}. Para calcular 18\% de 250, faça \texttt{0,18 x 250 = 45}. Um aumento de \texttt{p\%} multiplica por \texttt{1 + p/100}; uma redução multiplica por \texttt{1 - p/100}.

Percentuais sucessivos não devem ser somados automaticamente. Um produto que aumenta 20\% e depois diminui 20\% termina com fator \texttt{1,20 x 0,80 = 0,96}, ou seja, redução total de 4\%. Para descobrir o percentual de variação, use \texttt{(valor final - valor inicial) / valor inicial x 100}.

Em problemas de lucro, desconto e acréscimo, identifique a base: desconto de 15\% sobre 200 produz preço \texttt{200 x 0,85 = 170}. Se 170 é o preço após desconto de 15\%, o preço original não é 200; é \texttt{170/0,85 = 200}. Essa inversão aparece bastante em questões de preço e salário.

\subsection*{4. Equações e sistemas}

Uma equação é uma igualdade com incógnita. Preserve a igualdade fazendo a mesma operação nos dois membros. Na equação \texttt{3x - 7 = 11}, some 7 e divida por 3: \texttt{3x = 18}, portanto \texttt{x = 6}. mesma operação nos dois membros. Na equação \texttt{3x - 7 = 11}, some 7 e divida por 3: \texttt{3x = 18}, portanto \texttt{x = 6}.

Na equação quadrática \texttt{ax\textasciicircum{}2 + bx + c = 0}, com \texttt{a} diferente de zero, use \texttt{delta = b\textasciicircum{}2 - 4ac} e \texttt{x = (-b mais ou menos raiz de delta)/(2a)}. Se delta é positivo, há duas raízes reais; se é zero, uma raiz real dupla; se é negativo, não há raiz real. Quando possível, fatorar é mais rápido: \texttt{x\textasciicircum{}2 - 5x + 6 = 0} vira \texttt{(x - 2)(x - 3) = 0}, logo \texttt{x = 2} ou \texttt{x = 3}.

Em um sistema linear, elimine uma incógnita somando ou subtraindo equações. Para \texttt{x + y = 17} e \texttt{x - y = 5}, a soma dá \texttt{2x = 22}, então \texttt{x = 11} e \texttt{y = 6}. O método da substituição também é útil quando uma equação já está isolada. Confira os valores nas duas equações, principalmente em alternativas.

\subsection*{5. Produtos notáveis e fatoração}

Os produtos mais úteis são: \texttt{(a + b)\textasciicircum{}2 = a\textasciicircum{}2 + 2ab + b\textasciicircum{}2}; \texttt{(a - b)\textasciicircum{}2 = a\textasciicircum{}2 - 2ab + b\textasciicircum{}2}; e \texttt{(a + b)(a - b) = a\textasciicircum{}2 - b\textasciicircum{}2}. Também aparece \texttt{(a + b)\textasciicircum{}3 = a\textasciicircum{}3 + 3a\textasciicircum{}2b + 3ab\textasciicircum{}2 + b\textasciicircum{}3}.

Fatorar é escrever uma expressão como produto. Coloque o fator comum em evidência: \texttt{6x\textasciicircum{}2 + 9x = 3x(2x + 3)}. Na diferença de quadrados, \texttt{x\textasciicircum{}2 - 25 = (x - 5)(x + 5)}. No trinômio \texttt{x\textasciicircum{}2 + bx + c}, procure dois números cuja soma seja \texttt{b} e cujo produto seja \texttt{c}. A fatoração simplifica frações algébricas, resolve equações e revela zeros de funções.

\subsection*{6. Funções afim e quadrática}

Uma função associa cada entrada a uma saída. Na função afim \texttt{f(x) = ax + b}, \texttt{a} é a taxa de variação e \texttt{b} é o valor inicial, pois \texttt{f(0) = b}. Se \texttt{a > 0}, a função cresce; se \texttt{a < 0}, decresce. A raiz é \texttt{x = -b/a}, quando \texttt{a} é diferente de zero. Em um gráfico, \texttt{a} é a inclinação da reta.

Se dois pontos \texttt{(x1,y1)} e \texttt{(x2,y2)} são conhecidos, a taxa é \texttt{a = (y2 - y1)/(x2 - x1)}. Depois, use um dos pontos para achar \texttt{b}. Essa técnica modela tarifa fixa mais consumo, conversão de temperatura e custo de produção.

A função quadrática \texttt{f(x) = ax\textasciicircum{}2 + bx + c} tem gráfico parabólico. O vértice tem abscissa \texttt{xv = -b/(2a)} e ordenada \texttt{yv = -delta/(4a)}. Se \texttt{a > 0}, o vértice é mínimo; se \texttt{a < 0}, é máximo. As raízes, quando reais, indicam os pontos em que o gráfico cruza o eixo horizontal. Em problemas de área ou lucro, o vértice frequentemente representa o valor otimizado.

\subsection*{7. PA e PG}

Na progressão aritmética, a diferença entre termos é constante e recebe o nome de razão \texttt{r}. O termo geral é \texttt{an = a1 + (n - 1)r} e a soma dos \texttt{n} termos é \texttt{Sn = n(a1 + an)/2}. Se \texttt{a1 = 4} e \texttt{r = 3}, então \texttt{a5 = 4 + 4 x 3 = 16}.

Na progressão geométrica, cada termo é obtido multiplicando o anterior por uma razão \texttt{q}. O termo geral é \texttt{an = a1 x q\textasciicircum{}(n - 1)}. Para \texttt{q} diferente de 1, \texttt{Sn = a1(q\textasciicircum{}n - 1)/(q - 1)}. Se \texttt{|q| < 1}, a soma infinita é \texttt{S = a1/(1 - q)}. Em crescimento populacional, juros compostos e duplicações, procure PG; em aumentos constantes, procure PA.

\subsection*{8. Geometria plana e espacial}

Ângulos complementares somam 90 graus e suplementares somam 180 graus. A soma dos ângulos internos de um triângulo é 180 graus. Em um polígono de \texttt{n} lados, a soma é \texttt{(n - 2) x 180 graus}; cada ângulo interno de um polígono regular é essa soma dividida por \texttt{n}.

Áreas essenciais: retângulo \texttt{A = base x altura}; triângulo \texttt{A = base x altura/2}; paralelogramo \texttt{A = base x altura}; trapézio \texttt{A = (base maior + base menor) x altura/2}; círculo \texttt{A = pi x r\textasciicircum{}2}. Comprimento da circunferência é \texttt{2 x pi x r}. Use \texttt{pi} conforme a questão, normalmente 3,14 ou 22/7.

O teorema de Pitágoras vale no triângulo retângulo: \texttt{hipotenusa\textasciicircum{}2 = cateto1\textasciicircum{}2  + cateto2\textasciicircum{}2}. Em triângulos semelhantes, lados correspondentes são proporcionais e áreas variam com o quadrado da razão de semelhança.

Volumes: prisma e cilindro \texttt{V = área da base x altura}; pirâmide e cone \texttt{V = área da base x altura/3}; esfera \texttt{V = 4 x pi x r\textasciicircum{}3/3}. Áreas laterais e totais exigem distinguir a superfície lateral da base. Em sólidos semelhantes, comprimentos variam por \texttt{k}, áreas por \texttt{k\textasciicircum{}2} e volumes por \texttt{k\textasciicircum{}3}.

\subsection*{9. Trigonometria básica}

Em triângulo retângulo, para um ângulo agudo: \texttt{sen = cateto oposto/hipotenusa}, \texttt{cos = cateto adjacente/hipotenusa} e \texttt{tg = cateto oposto/cateto adjacente}. Os valores notáveis são: 30 graus, \texttt{sen = 1/2}, \texttt{cos = raiz de 3/2}, \texttt{tg = raiz de 3/3}; 45 graus, \texttt{sen = cos = raiz de 2/2}, \texttt{tg = 1}; 60 graus, \texttt{sen = raiz de 3/2}, \texttt{cos = 1/2}, \texttt{tg = raiz de 3}.

A identidade fundamental é \texttt{sen\textasciicircum{}2(x) + cos\textasciicircum{}2(x) = 1}. Para alturas e distâncias, faça um desenho, indique o ângulo e identifique os lados antes de escolher a razão. Se o ângulo está no solo e a altura é o lado oposto, a tangente costuma ser a mais direta.

\subsection*{10. Análise combinatória e probabilidade}

O princípio multiplicativo diz que, se uma escolha tem \texttt{m} opções e a seguinte tem \texttt{n}, o total é \texttt{m x n}. Permutação de \texttt{n} objetos distintos é \texttt{n!}. Arranjos de \texttt{n} elementos tomados \texttt{p} a \texttt{p} são \texttt{A(n,p) = n!/(n-p)!}. Combinações são \texttt{C(n,p) = n!/[p!(n-p)!]}, usadas quando a ordem não importa. Sempre pergunte: trocar a ordem produz resultado diferente?

Probabilidade em espaço equiprovável é \texttt{P(A) = casos favoráveis/casos possíveis}. O complemento é \texttt{P(não A) = 1 - P(A)}. Para eventos independentes, multiplique probabilidades. Para eventos mutuamente exclusivos, some-as. Em situações condicionais, \texttt{P(A dado B) = P(A e B)/P(B)}. Diagramas de árvore ajudam quando há retiradas sucessivas sem reposição, pois o total muda após cada retirada.

\subsection*{11. Estatística}

A média aritmética é a soma dos valores dividida pela quantidade. Se os dados  têm pesos, use média ponderada: \texttt{soma(valor x peso)/soma dos pesos}. A mediana é o termo central após ordenar; com quantidade par, é a média dos dois centrais. A moda é o valor mais frequente.

Amplitude é \texttt{maior valor - menor valor}. A variância mede a dispersão em torno da média e o desvio padrão é sua raiz quadrada. Para questões escolares, a leitura de gráficos exige conferir escala, unidade e se o eixo começa em zero. Uma média alta não significa que todos os valores são altos; compare também mediana e amplitude.

\subsection*{12. Matemática financeira}

Em juros simples, os juros incidem sempre sobre o capital inicial: \texttt{J = C x i x t} e \texttt{M = C + J = C(1 + i x t)}. Em juros compostos, há juros sobre juros: \texttt{M = C(1 + i)\textasciicircum{}t}, com \texttt{J = M - C}. A taxa e o tempo devem estar na mesma unidade. Uma taxa mensal não pode ser aplicada diretamente a um prazo anual sem conversão.

Desconto simples por fora costuma usar \texttt{D = N x i x t}, em que \texttt{N} é o valor nominal, e o valor atual é \texttt{A = N - D}. Em uma sequência de descontos, use os fatores multiplicativos: desconto de 10\% seguido de 5\% deixa \texttt{0,90 x 0,95 = 0,855} do valor inicial, equivalendo a desconto total de 14,5\%.

\subsubsection*{Aplicação e conferência}

Nos problemas de operações, faça uma estimativa antes do cálculo exato. Se \texttt{398 x 0,24} aparece, o resultado deve ficar próximo de 100, não de 1.000. Para decimais, retire temporariamente as vírgulas, efetue a operação com inteiros e reposicione a vírgula contando as casas decimais dos fatores. Em expressões com frações, a barra deve ser entendida como parênteses: em \texttt{1/(2 + 3)}, o denominador é todo \texttt{2 + 3}, e não apenas o número 2.

Na regra de três, não coloque grandezas diferentes na mesma coluna sem indicar o que está sendo comparado. Se uma receita para 4 pessoas usa 300 gramas de arroz, a relação \texttt{300/4} tem unidade gramas por pessoa. Para 10 pessoas, \texttt{300 x 10/4 = 750} gramas. Já em produtividade, mais trabalhadores diminuem o tempo, portanto o produto trabalhadores x tempo deve ser preservado. Uma tabela com setas para cima e para baixo ajuda a evitar a inversão.

Em percentuais, a palavra "de" indica multiplicação. Assim, 30\% de 80 é 24, mas aumentar 80 em 30\% produz 104. Em um problema de população, saldo ou estoque, use o valor inicial como base de cada alteração. Se a base mudar, o percentual equivalente também muda. Essa é a razão de dois descontos de 10\% representarem redução de 19\%, e não de 20\%: \texttt{0,90 x 0,90 = 0,81}.

Ao montar uma equação, defina a incógnita com unidade. Se \texttt{x} é a idade atual, "daqui a 5 anos" vira \texttt{x + 5}, e "há 3 anos" vira \texttt{x - 3}. Se \texttt{x} é uma quantidade de objetos, imponha que \texttt{x} seja inteiro e não negativo. Em equações com denominadores, exclua valores que zeram o denominador antes de multiplicar os dois membros. Em problemas de velocidade, use \texttt{distância = velocidade x tempo} e converta km/h para m/s multiplicando por \texttt{5/18}.

Para funções, uma tabela com alguns pares \texttt{(x, f(x))} é uma forma de verificar o modelo. Na afim, diferenças iguais em intervalos iguais indicam crescimento linear. Na quadrática, diferenças de primeira ordem variam, mas diferenças de segunda ordem são constantes. O domínio pode ser restringido pelo contexto: tempo não é negativo, número de pessoas é inteiro e dimensões são positivas. Não aceite uma raiz algébrica que viola essas condições.

Em progressões, identifique o índice inicial. Se a sequência começa em \texttt{a0}, o termo geral pode aparecer como \texttt{an = a0 + n x r} na PA ou \texttt{an = a0 x q\textasciicircum{}n} na PG. Compare a fórmula com os primeiros termos para não deslocar o índice. Em soma de PA, emparelhe primeiro e último termo: todos os pares têm a mesma soma. Em PG, conte quantas multiplicações separam o primeiro termo do termo solicitado.

Na geometria, decomponha figuras compostas em partes conhecidas e some áreas; para regiões vazias, subtraia a área interna da externa. Perímetro mede contorno e área mede superfície, por isso não se somam diretamente. Ao ampliar um desenho, uma escala \texttt{1:50} significa que 1 unidade no desenho representa 50 na realidade; para área, o fator é \texttt{50\textasciicircum{}2}, e para volume, \texttt{50\textasciicircum{}3}. Em uma figura espacial, conte as dimensões antes de escolher a fórmula.

Na trigonometria, os valores notáveis podem ser reconstruídos com os triângulos 30-60-90 e 45-45-90. Não confunda seno com cosseno: ambos usam a hipotenusa, mas trocam os catetos. Se o problema fornece dois lados e pede um ângulo, use a razão conhecida e, se necessário, a tabela ou a calculadora autorizada. Em distâncias inacessíveis, a tangente relaciona altura e distância horizontal sem exigir a hipotenusa.

Na análise combinatória, conte primeiro sem restrições e depois trate as restrições. Para objetos juntos, considere o bloco; para posições proibidas, subtraia os casos inválidos. Se há repetição, a permutação de \texttt{n} objetos com repetições de quantidades \texttt{a}, \texttt{b} e \texttt{c} é \texttt{n!/(a! x b! x c!)}. Em probabilidade, uma contagem bem feita vem antes da divisão. Quando o evento pedido é "ao menos um", frequentemente é mais rápido calcular o complemento "nenhum".

Em estatística, ordenar os dados é obrigatório para achar mediana e quartis. Se todos os valores recebem a mesma constante, a média e a mediana aumentam por essa constante, mas a amplitude não muda. Se todos são multiplicados por um fator positivo, média, mediana e amplitude são multiplicadas pelo fator. Em gráficos, compare variações absolutas e relativas: passar de 20 para 30 é mais 10 unidades, mas crescimento relativo de 50\%.

Em finanças, transforme 2\% em \texttt{0,02}, nunca em \texttt{2}. Se o prazo é 18 meses e a taxa é anual simples, use \texttt{t = 18/12 = 1,5 ano}, desde que a convenção da questão seja proporcional. Em compostos, taxa efetiva anual e taxa mensal não são convertidas apenas multiplicando por 12: uma taxa mensal \texttt{i} corresponde a \texttt{(1 + i)\textasciicircum{}12 - 1} ao ano. Em parcelas, não compare somente o valor mensal; compare o montante total e identifique se há entrada, desconto ou juros.

\section*{Como a banca cobra}

\begin{itemize}
  \item Em cálculos longos, a alternativa costuma denunciar erro de sinal ou de unidade; estime o resultado antes de finalizar.
  \item Frações de uma quantidade são multiplicações, mas a soma de frações exige denominador comum.
  \item Regra direta acompanha o crescimento; regra inversa conserva o produto.
  \item Percentuais sucessivos são multiplicados, não somados sem análise.
  \item Em equações, a solução deve ser testada no enunciado e respeitar domínio, como comprimento positivo.
  \item Na função, diferencie valor inicial, taxa de variação, raiz e vértice.
  \item Em geometria, área usa unidade ao quadrado e volume usa unidade ao cubo.
  \item Em combinatória, decidir se a ordem importa é mais importante que decorar fórmulas.
  \item Probabilidade sem reposição altera o espaço de casos possíveis.
  \item Na financeira, juros simples são lineares e compostos são exponenciais.
\end{itemize}

\subsection*{Exemplos resolvidos}

\subsubsection*{Exemplo 1 — porcentagem e preço}

Uma inscrição de 240 reais recebe desconto de 12\% para pagamento antecipado. O desconto é \texttt{0,12 x 240 = 28,80}. Logo, o preço pago é \texttt{240 - 28,80 = 211,20} reais. Se o preço final fosse dado e o inicial fosse pedido, dividiríamos pelo fator \texttt{0,88}, em vez de subtrair 12\% do valor final.

\subsubsection*{Exemplo 2 — sistema aplicado}

Em uma papelaria, 2 cadernos e 3 canetas custam 31 reais, enquanto 1 caderno e 2 canetas custam 19 reais. Seja \texttt{c} o preço do caderno e \texttt{p} o da caneta: \texttt{2c + 3p = 31} e \texttt{c + 2p = 19}. Multiplicando a segunda por 2, temos \texttt{2c + 4p = 38}. Subtraindo a primeira, \texttt{p = 7}; então \texttt{c + 14 = 19}, logo \texttt{c = 5}. A substituição final confirma os dois preços.

\subsubsection*{Exemplo 3 — função quadrática}

Um terreno retangular tem perímetro 40 metros. Se uma dimensão é \texttt{x}, a outra é \texttt{20 - x}, e a área é \texttt{A(x) = x(20 - x) = -x\textasciicircum{}2 + 20x}. Como \texttt{a = -1}, há máximo no vértice: \texttt{xv = -20/(2 x -1) = 10}. A outra dimensão também é 10 e a área máxima é \texttt{100 m2}. O domínio físico é \texttt{0 < x < 20}.

\subsubsection*{Exemplo 4 — contagem e probabilidade}

De um grupo de 5 candidatos, serão escolhidos 2 para uma comissão. Como a ordem não importa, o total é \texttt{C(5,2) = 5!/(2! x 3!) = 10}. Se apenas 3 dos candidatos têm curso de primeiros socorros, a chance de os dois escolhidos terem esse curso é \texttt{C(3,2)/C(5,2) = 3/10}, ou 30\%.

\subsubsection*{Exemplo 5 — juros compostos}

Um capital de 1.000 reais rende 10\% ao mês por 2 meses. O montante é \texttt{M = 1000 x (1,10)\textasciicircum{}2 = 1000 x 1,21 = 1210}. Portanto, os juros são \texttt{1210 - 1000 = 210} reais. Em juros simples, no mesmo período, seriam \texttt{1000 x 0,10 x 2 = 200} reais.

\section*{Quadro-resumo}

\begin{tabularx}{\linewidth}{llX}
\toprule
 \textbf{Tópico} & \textbf{Relação principal} & \textbf{Atenção de prova} \\
\midrule
 Frações & somar com denominador comum; dividir pelo inverso & simplificar e conservar sinais \\
 Proporção & \texttt{a/b = c/d} implica \texttt{a x d = b x c} & conferir grandezas diretas ou inversas \\
 Porcentagem & aumento: \texttt{valor x (1 + taxa)} & percentuais sucessivos multiplicam fatores \\
 Equação quadrática & \texttt{delta = b\textasciicircum{}2 - 4ac} & respeitar domínio e interpretar raízes \\
 Função afim & \texttt{f(x) = ax + b} & \texttt{a} é a inclinação e \texttt{b} é o valor inicial \\
 PA & \texttt{an = a1 + (n - 1)r} & diferença constante \\
 PG & \texttt{an = a1 x q\textasciicircum{}(n - 1)} & razão multiplicativa \\
 Triângulo & \texttt{A = base x altura/2} & altura é perpendicular à base \\
 Sólidos & \texttt{V = área da base x altura} em prismas & unidade cúbica \\
 Pitágoras & \texttt{h\textasciicircum{}2 = c1\textasciicircum{}2 + c2\textasciicircum{}2} & somente triângulo retângulo \\
 Combinação & \texttt{C(n,p) = n!/[p!(n-p)!]} & ordem não importa \\
 Probabilidade & favoráveis sobre possíveis & sem reposição muda os denominadores \\
 Média & soma dividida pela quantidade & média ponderada usa pesos \\
 Juros simples & \texttt{M = C(1 + i x t)} & base é sempre o capital inicial \\
 Juros compostos & \texttt{M = C(1 + i)\textasciicircum{}t} & alinhar taxa e período \\
\bottomrule
\end{tabularx}


\section*{Questões de fixação}

\subsection*{Questão 1}

Uma equipe percorreu \texttt{3/8} de uma trilha de 24 km pela manhã e \texttt{1/4} da trilha à tarde. Qual distância ainda faltou?

A) 6 km \\ B) 8 km \\ C) 9 km \\ D) 11 km \\ E) 15 km

\begin{alerta}
\textbf{Gabarito: C.} As partes percorridas somam \texttt{3/8 + 1/4 = 3/8 + 2/8 = 5/8}. Isso corresponde a \texttt{24 x 5/8 = 15} km. Restam \texttt{24 - 15 = 9} km.
\end{alerta}

\subsection*{Questão 2}

Se 8 máquinas produzem uma encomenda em 15 horas, quantas horas serão necessárias para 10 máquinas iguais, no mesmo ritmo?

A) 10 \\ B) 12 \\ C) 13 \\ D) 18 \\ E) 20

\begin{alerta}
\textbf{Gabarito: B.} É uma relação inversa: máquinas x horas é constante. \texttt{8 x 15 = 10 x t}, então \texttt{t = 120/10 = 12} horas.
\end{alerta}

\subsection*{Questão 3}

Um produto de 320 reais sofre aumento de 25\% e depois desconto de 20\% sobre o novo preço. O preço final é:

A) 300 reais \\ B) 320 reais \\ C) 325 reais \\ D) 336 reais \\ E) 360 reais

\begin{alerta}
\textbf{Gabarito: B.} O aumento gera \texttt{320 x 1,25 = 400}. O desconto gera \texttt{400 x 0,80 = 320}. Os percentuais se anulam apenas porque os fatores são inversos neste caso.
\end{alerta}

\subsection*{Questão 4}

O sistema \texttt{2x + y = 13} e \texttt{x - y = 2} tem como solução:

A) \texttt{(3, 4)} \\ B) \texttt{(4, 5)} \\ C) \texttt{(5, 3)} \\ D) \texttt{(6, 1)} \\ E) \texttt{(7, -1)}

\begin{alerta}
\textbf{Gabarito: C.} Somando as equações, \texttt{3x = 15}, logo \texttt{x = 5}. Na segunda, \texttt{5 - y = 2}, portanto \texttt{y = 3}.
\end{alerta}

\subsection*{Questão 5}

O quinto termo da PA \texttt{7, 11, 15, ...} e a soma dos cinco primeiros termos são,

A) 19 e 55 \\ B) 23 e 75 \\ C) 23 e 85 \\ D) 27 e 85 \\ E) 27 e 95

\begin{alerta}
\textbf{Gabarito: B.} A razão é 4. \texttt{a5 = 7 + 4 x 4 = 23}. A soma é \texttt{S5 = 5(7 + 23)/2 = 5 x 30/2 = 75}, portanto a alternativa correta é B.
\end{alerta}

\subsection*{Questão 6}

Um retângulo tem diagonal 13 cm e um dos lados mede 5 cm. Sua área é:

A) 30 cm2 \\ B) 40 cm2 \\ C) 60 cm2 \\ D) 65 cm2 \\ E) 78 cm2

\begin{alerta}
\textbf{Gabarito: C.} Pelo teorema de Pitágoras, o outro lado satisfaz \texttt{13\textasciicircum{}2 = 5\textasciicircum{}2 + x\textasciicircum{}2}; então \texttt{x\textasciicircum{}2 = 169 - 25 = 144} e \texttt{x = 12}. A área é \texttt{5 x 12 = 60 cm2}.
\end{alerta}

\subsection*{Questão 7}

Uma urna tem 4 bolas brancas e 6 pretas. Retirando uma bola ao acaso, a probabilidade de ela não ser preta é:

A) 20\% \\ B) 30\% \\ C) 40\% \\ D) 60\% \\ E) 70\%

\begin{alerta}
\textbf{Gabarito: C.} Não ser preta significa ser branca. Há 4 casos favoráveis em 10 possíveis: \texttt{4/10 = 0,4 = 40\%}.
\end{alerta}

\subsection*{Questão 8}

As notas 6, 8 e 10 têm pesos 1, 2 e 2, respectivamente. A média ponderada é:

A) 7,2 \\ B) 7,6 \\ C) 8,0 \\ D) 8,2 \\ E) 8,4

\begin{alerta}
\textbf{Gabarito: E.} A soma ponderada é \texttt{6 x 1 + 8 x 2 + 10 x 2 = 42}. A soma dos pesos é 5. Assim, a média é \texttt{42/5 = 8,4}. A média simples não deve ser usada porque as notas têm pesos diferentes.
\end{alerta}

\subsection*{Questão 9}

Um capital de 2.000 reais é aplicado a juros simples de 3\% ao mês durante 4 meses. O montante obtido é:

A) 2.060 reais \\ B) 2.120 reais \\ C) 2.180 reais \\ D) 2.240 reais \\ E) 2.360 reais

\begin{alerta}
\textbf{Gabarito: D.} Os juros são \texttt{J = 2000 x 0,03 x 4 = 240}. Logo, \texttt{M = 2000 + 240 = 2240} reais.
\end{alerta}

\subsection*{Questão 10}

Quantos números de três algarismos distintos podem ser formados com os dígitos 1, 2, 3, 4 e 5?

A) 10 \\ B) 15 \\ C) 20 \\ D) 30 \\ E) 60

\begin{alerta}
\textbf{Gabarito: E.} Para a centena há 5 escolhas, para a dezena restam 4 e para  a unidade restam 3. Pelo princípio multiplicativo, \texttt{5 x 4 x 3 = 60}.
\end{alerta}

\section*{Revisão final}

\begin{enumerate}
  \item Resolva operações respeitando a ordem e reduza frações antes de comparar.
  \item Em proporções, verifique se a relação é direta ou inversa e uniformize as unidades.
  \item Converta porcentagem em fator: \texttt{1 + taxa} para aumento e \texttt{1 - taxa} para desconto.
  \item Em sistemas, elimine uma incógnita e confira a solução nas equações originais.
  \item Use produtos notáveis para expandir e fatoração para simplificar ou encontrar raízes.
  \item Na função afim, interprete a inclinação; na quadrática, procure raízes e vértice.
  \item Diferencie PA, com diferença constante, de PG, com razão multiplicativa.
  \item Desenhe figuras, use Pitágoras apenas em triângulos retângulos e controle unidades de área e volume.
  \item Em contagem, decida se a ordem importa; em probabilidade, defina casos possíveis e favoráveis.
  \item Leia média, mediana e dispersão juntas e alinhe taxa e tempo em juros.
\end{enumerate}

Na prova, não abandone a estimativa. Um comprimento negativo, uma probabilidade maior que 1 ou uma área incompatível com o desenho indicam erro de modelagem. A sequência mais segura é: ler, destacar dados, escrever a relação, calcular, testar e marcar a alternativa. Em questões de múltipla escolha, substitua as alternativas quando a expressão permitir uma verificação rápida; em questões discursivas, deixe claras as unidades e as etapas intermediárias. Se o resultado for decimal periódico, mantenha a fração até o fim para evitar arredondamentos. Quando duas alternativas parecem próximas, volte ao enunciado e verifique se a pergunta pede valor, variação, porcentagem, área, perímetro ou probabilidade. Essa rotina reduz erros de leitura e mantém o cálculo objetivo, tanto em provas de concurso quanto no ENEM.

\end{document}
